\section{Conclusion}

A VA estimate is one point in the plane, so its uncertainty should be one region
in the plane. Calibrating valence and arousal separately, as is common, leaves
the guarantee on the wrong event and under-covers the VA point by about
four to five points at nominal $90\%$. One max-normalised score fixes this with a
finite-sample joint guarantee and smaller regions than the Bonferroni correction.
It costs a single line of code, and we recommend it as the default.

The harder finding is where guarantees end. An expression absent from the
calibration data loses coverage under every method we tried, and none of the
uncertainty estimators reliably notices. Conformal methods cannot fix this by
construction, so a deployed affect estimator needs either calibration data that
spans the expressions it will meet or an uncertainty signal that actually rises
on novel ones. We found neither in the standard toolkit, and the reason is
instructive: contempt is not an unusual input but a familiar-looking face with an
unfamiliar meaning, so no signal computed from the input alone can be expected
to flag it.

\paragraph{Limitations.} We study still images from one dataset family and two
backbones (the second with three seeds). A
cross-dataset shift (AFEW-VA, Aff-Wild2) and annotator-level labels would test
whether the uncertainty tracks human disagreement. Our unseen expression has only
$70$ examples, so its coverage estimates carry wide intervals (we report them).
The categorical labels used for the per-expression audit are themselves noisy
annotations of a continuous state.
